%KerMLExecution.tex

\section{Overview}

\comment{The method of using \texttt{\#atom} to define instances as KerML text was at first confusing, and then unsatisfying.  The  \texttt{\#atom} really doesn't say anything.
An ``instance'' is not created--just repeat (mostly) of the text of the model which is supposed to be executed.}

\section{Modeling Instances}
\paragraph{Classifier Specialization}

The first example shows classifier specialization and multiplicity.

\begin{lstlisting}
classifier Vehicle;
classifier Bicycle specializes Vehicle;
classifier MyBike [1] specializes Bicycle;
classifier YourBike [1] specializes Bicycle disjoint from MyBike;
classifier OurBikes [2] unions MyBike, YourBike;
\end{lstlisting}

By \ref{eq:df-classifier}:
\begin{gather}
 (\texttt{classifier Vehicle;} \equiv \texttt{Vehicle}\in V_C~)
\end{gather}
 
By \ref{eq:df-rep} \texttt{Vehicle} is represented as a triple:
\begin{gather}
\langle~\textbf{classifier},\texttt{"Vehicle"},\{x~\vert ~(~ x \in \texttt{Vehicle} \wedge \texttt{Vehicle}\in V_C~)~\}~\rangle
\label{eq:vehicle}
\end{gather}
 
By \ref{eq:df-classifier} and \ref{eq:df-specializes} :
 \begin{gather}
 (\texttt{classifier Bicycle specializes Vehicle;} \equiv \notag \\
  (~\texttt{Bicycle}\in V_C~\wedge~ \texttt{Bicycle}\subseteq\texttt{Vehicle}~)~)
\end{gather}
 
By \ref{eq:df-rep} \texttt{Bicycle} is represented as a triple:
\begin{gather}
\langle~\textbf{classifier},\texttt{"Bicycle"},\{x~\vert ~(~ x \in \texttt{Bicycle} ~\wedge~ \texttt{Bicycle}\in V_C~\wedge~ \texttt{Bicycle}\subseteq\texttt{Vehicle}~)~\}~\rangle
\label{eq:bicycle}
\end{gather}

By \ref{eq:df-classifier}, \ref{eq:df-specializes}, and \ref{eq:df-typemultiplicity} :
\begin{gather}
(\texttt{classifier MyBike [1] specializes Bicycle;} \equiv \notag \\
(~\texttt{MyBike}\in V_C~\wedge~ \texttt{MyBike}\subseteq\texttt{Bicycle}~\wedge \#\texttt{MyBike}=1~)~)
\end{gather}
 
By \ref{eq:df-rep} \texttt{MyBike} is represented as a triple:
\begin{gather}
\langle~\textbf{classifier},\texttt{"MyBike"},\{x~\vert ~(~ x \in \texttt{MyBike}~ \wedge \notag \\
 ~~~\texttt{MyBike}\in V_C~\wedge~ \texttt{MyBike}\subseteq\texttt{Bicycle}~\wedge \#\texttt{MyBike}=1~)~\}~\rangle
\label{eq:mybike}
\end{gather}

By \ref{eq:df-classifier}, \ref{eq:df-specializes}, \ref{eq:df-typemultiplicity}, and \ref{eq:df-disjoint} :
\begin{gather}
(\texttt{classifier YourBike [1] specializes Bicycle disjoint from MyBike;} \equiv \notag \\
(~\texttt{YourBike}\in V_C~\wedge~ \texttt{YourBike}\subseteq\texttt{Bicycle}~\wedge \#\texttt{YourBike}=1~\wedge~\texttt{YourBike}\cap\texttt{MyBike}=\emptyset)~)
\end{gather}
 
By \ref{eq:df-rep} \texttt{YourBike} is represented as a triple:
\begin{gather}
\langle~\textbf{classifier},\texttt{"YourBike"},\{x~\vert ~(~ x \in \texttt{YourBike}~\wedge~\texttt{YourBike}\in V_C~ \wedge \notag \\
 ~~~ \texttt{YourBike}\subseteq\texttt{Bicycle}~\wedge \#\texttt{YourBike}=1~\wedge~\texttt{YourBike}\cap\texttt{MyBike}=\emptyset~)~\}~\rangle
\label{eq:yourbike}
\end{gather}

By \ref{eq:df-unions}:
\begin{gather}
(\texttt{classifier OurBikes unions MyBike, YourBike;} \equiv \notag \\
(~\texttt{OurBikes}\in V_C~\wedge~ \texttt{OurBikes}=\texttt{MyBike}\cup\texttt{YourBike}~\wedge~\#\texttt{OurBikes}=2~)~)
\end{gather}
By \ref{eq:df-rep} \texttt{OurBikes} is represented as a triple:
\begin{gather}
\langle~\textbf{classifier},\texttt{"OurBikes"},\{x~\vert ~(~ x \in \texttt{OurBikes}~\wedge~\texttt{OurBikes}\in V_C~ \wedge \notag \\
 ~~~  \texttt{OurBikes}=\texttt{MyBike}\cup\texttt{YourBike}~\wedge~\#\texttt{OurBikes}=2~)~\}~\rangle
\label{eq:ourbikes}
\end{gather}

So, the design, $D$ a set of triples \ref{eq:vehicle}, \ref{eq:bicycle}, \ref{eq:mybike}, \ref{eq:yourbike}, \ref{eq:ourbikes}, stored in index M.

This is a fine example how ZFC-classes get defined for KerML elements, but it's not really execution.  By \ref{eq:df-model} execution means constructing $\mathfrak{M}$ by selecting one element, $E$, to be instantiated.  Constructing $\mathfrak{M}$ requires creating the interpretation $\mathbf{I}\lsem d,\tau\rsem$.  This gets more challenging when the system interacts with its environment at its system boundary, $P(\tau)$.
In this example, nothing changes over time, so the interpretation is the same for all $\tau\in\textbf{T}$.  

\section{Instantiation Procedure}

\subsection{Overview}

Creating a model $\mathfrak{M}$ means creating an interpretation, $\mathbf{I}\lsem E,\tau\rsem$, for the selected element of the design $E\in D$, for all $\tau\in\mathbf{T}$.
Fortunately, we can model a discrete, finite, set of instants where something changes, and assume that the interpretation stays the same until the next instant where something changes.

Suppose there are $k$ instants between $0$ and $\mathsf{\mc{now}}$ at which something changes.  Let the set of change instants be $X_\tau$.

\pn Not every design element has a value that can meaningfully change over time.  A plain classifier
like \texttt{Vehicle} (\ref{eq:vehicle}) denotes a ZFC-class; $\mathbf{I}\lsem\texttt{Vehicle},\tau\rsem$
\emph{selects} a member of that class, it doesn't return a value that itself varies.  By \ref{eq:df-rep},
$k(\texttt{Vehicle})=\textbf{classifier}$, while a \textbf{feature}'s interpretation is exactly the kind
of time-varying value change/no-change reasoning needs.  Let $\mathit{TVal}\subseteq D$ be the subset of
the design whose kind is \textbf{feature}:

\begin{restatable}{definition}{dftval}~\blTitle{df-tval}~ Define the subset of the design with a time-varying value.
\begin{gather}
\mathit{TVal}~=~\{~d\in D~\vert~\exists id,C~(~d~=~\langle~\textbf{feature},id,C~\rangle~)~\}
\label{eq:df-tval}\tag{df-tval}
\end{gather}
\end{restatable}

\pn \texttt{df-tval} is book-original, with no \texttt{SFS.mm} formalization of its own -- it names, at
the book level, exactly the domain \texttt{SFS.lean}'s \texttt{TVal~$\alpha$}
(\texttt{:=~Time~$\to$~$\alpha$}) already covers: a value that is literally a function of time, as
opposed to a classifier's class-membership selection.

\begin{restatable}{definition}{dfchange}~\blTitle{df-change}~ Define set of time instants at which something changes.
\begin{gather}
(~X_\tau~=~\{~\tau_1,\tau_2,\ldots,\tau_k~\}~\wedge~X_\tau~\subset~\mathbf{T}~\wedge \notag \\
~\forall \tau \in \mathbf{T}~\vert~((\exists d \in \mathit{TVal}~\exists \tau_i\prec\tau~\vert~ \notag \\
~~~(\forall \tau'~\vert~\tau_i\prec\tau'\prec\tau~\rightarrow~\mathbf{I}\lsem d,\tau'\rsem=\mathbf{I}\lsem d,\tau_i\rsem)~\wedge~\mathbf{I}\lsem d,\tau_i\rsem\neq\mathbf{I}\lsem d,\tau\rsem)~\leftrightarrow~\tau \in X_\tau )~\wedge \notag \\
~\forall i,j\in(1[,]k)~\vert~(~i<j~\rightarrow \tau_i\prec\tau_j~)~) \label{eq:df-change}\tag{df-change}
\end{gather}
\end{restatable}

\pn \texttt{df-change} is stated directly in terms of $\mathbf{I}\lsem d,\tau\rsem$ rather than the
\texttt{changed} predicate used elsewhere in \texttt{SFS.mm} and \texttt{SFS.lean}: that \texttt{changed}
(\texttt{SFS.mm} \texttt{cchanged} and \texttt{df-bl.changed}; \texttt{SFS.lean}'s
\texttt{changed~:~}\texttt{Oc\-currence~$\to$~Time~$\to$~Time}) is domain-restricted to \texttt{Occurrence}, while
$\mathit{TVal}$ (\ref{eq:df-tval}) plays the same domain-restricting role in general, matching
\texttt{SFS.lean}'s own \texttt{TVal~$\alpha$} rather than one specific to \texttt{Occurrence}. $X_\tau$
no longer forces $0$ or $\mathsf{now}$ to be members: $0\prec\tau_1$ is automatic (the
$\exists\tau_i\prec\tau$ hypothesis can never be satisfied at $\tau=0$, since nothing in $\mathbf{T}$
precedes it), and $\tau_k$ may legitimately equal $\mathsf{now}$ if something changes at the last
observed instant -- both were artificial exclusions of the earlier boundary-forcing form. $\tau$ is a
change instant exactly when some $d\in\mathit{TVal}$ held a constant interpretation on an interval right
up to $\tau$ (from some earlier point $\tau_i$) and that constant value differs from
$\mathbf{I}\lsem d,\tau\rsem$ itself -- matching \texttt{df-nochange}'s own constant-on-an-interval
vocabulary below, without presupposing $X_\tau$'s own enumeration. 

\pn Stated this way, a \texttt{Lean}
counterpart of \texttt{df-change} can be \emph{proved}, not just asserted, from \texttt{SFS.lean}'s
\texttt{finitePartition} axiom: its contrapositive (no shared tick in $(t_1,t_2]$ implies equal
interpretation at $t_1,t_2$) is exactly what forces every $\tau$ satisfying this jump condition to be one
of \texttt{finitePartition}'s own shared \texttt{breaks}, giving $X_\tau\subseteq\{\texttt{breaks}(i)\}$
-- hence finite -- as a theorem, rather than assuming finiteness by writing $X_\tau$ as a literal finite
set. (The converse doesn't hold -- not every shared tick need be an actual change -- which is exactly why
$X_\tau$ must be \emph{defined} by the jump condition, as above, rather than \emph{as} the break set
itself.)

\begin{restatable}{leandefinition}{ChangeSetLean}~\leanTitle{changeSet}~ \texttt{df-change}'s $X_\tau$,
restated for a single $d$ over the concrete shared ticks (\texttt{ticks}, extracted from
\texttt{finitePartition} via \texttt{Classical.choose}) rather than an open existential over
\texttt{Time}: a genuine change can only ever be \emph{at} a tick (\texttt{ticks\_bucket}), never
strictly between two of them, so restricting to ticks loses no generality.

\texttt{\textcolor{leancolor}{def} \textcolor{leanyellow}{changeSet} \{$\alpha$ : Type\} (d : TVal $\alpha$) : Set Time :=}

\texttt{~~\{$\tau$ : Time $\vert$ $\exists$ i : Fin (tickCount + 1), ticks i = $\tau$ $\wedge$}

\texttt{~~~~$\exists$ h : 0 < (i : $\mathbb{N}$), interpValAt d (ticks i) $\neq$}

\texttt{~~~~~~interpValAt d (ticks $\langle$(i : $\mathbb{N}$) - 1, by omega$\rangle$)\}}
\label{changeSet}
\end{restatable}

\begin{restatable}{leantheorem}{ChangeSetFiniteLean}~\leanTitle{changeSet\_finite}~ \texttt{df-change}'s
finiteness claim, proved rather than assumed: \texttt{changeSet d} is, by construction
(\texttt{changeSet\_subset\_ticks}), a subset of \texttt{Set.range ticks} -- finite because
\texttt{Fin (tickCount+1)} is. \texttt{finitePartition} does the real work one level down, inside
\texttt{ticks} itself (extracted from it via \texttt{Classical.choose}), not in this proof directly.

\texttt{\textcolor{leancolor}{theorem} \textcolor{leanyellow}{changeSet\_finite} \{$\alpha$ : Type\} (d : TVal $\alpha$) : (changeSet d).Finite :=}

\texttt{~~(Set.finite\_range ticks).subset (changeSet\_subset\_ticks d)}
\label{changeSet_finite}
\end{restatable}

When making $\mathbf{I}\lsem E,\tau\rsem$ the number of instants needed may not be known.  Furthermore, in may be necessary to introduce instants between others.
Therefore, instants will be named $\tau_a,\tau_b,\ldots$.  If an instant is needed between $\tau_a$ and $\tau_b$, call it $\tau_{a.a}$, and so on.

Between instants where something changes, it stays the same.  Let $X_\tau$ be the set of time instants at which something changes.
\begin{restatable}{definition}{dfnochange}~\blTitle{df-nochange}~ Interpretation stay the same between changes.
\begin{gather}
X_\tau~=~\{~\tau_1,\tau_2,\ldots,\tau_k~\}~\wedge \notag \\
\forall d\in \mathit{TVal}~\vert~(~\forall \tau\in (0(,) \tau_1)~\vert~\mathbf{I}\lsem d,\tau\rsem=\mathbf{I}\lsem d,0 \rsem ~\wedge \notag \\
~~~\forall \tau\in (\tau_k(,) \mathsf{\mc{now}})~\vert~\mathbf{I}\lsem d,\tau\rsem=\mathbf{I}\lsem d,\tau_k \rsem ~\wedge \notag \\
~~~\forall i\in(1[,]k\minus 1)~\vert~\forall \tau\in (\tau_i(,) \tau_{i+1})~\vert~\mathbf{I}\lsem d,\tau\rsem=\mathbf{I}\lsem d,\tau_i \rsem )\label{eq:df-nochange}\tag{df-nochange}
\end{gather}
\end{restatable}

\begin{restatable}{leantheorem}{ChangeSetConstantLean}~\leanTitle{changeSet\_constant}~ \texttt{df-nochange}:
if no element of \texttt{changeSet d} falls in $(t_1,t_2]$, $d$'s value is the same at $t_1$ and $t_2$ --
sharper than \texttt{SFS.lean}'s own \texttt{ticks\_constant} (which needs no tick at all in range, not
just no \emph{changing} tick: a "wasted" tick with no real jump is allowed inside). Proved by chaining,
tick index by tick index, from $t_1$'s largest-tick-below anchor up to $t_2$'s: at each step, the next
tick isn't a change (excluded by \texttt{hno}), so its value matches the one before it
(\texttt{changeSet\_step}, the contrapositive of \texttt{changeSet\_iff}). Uses
\texttt{finitePartition} both indirectly (anchoring $t_1$/$t_2$ to their nearest ticks) and directly (via
\texttt{changeSet}'s own dependence on \texttt{ticks}) -- confirmed via \texttt{\#print axioms}: no axiom
beyond the project's established baseline (\texttt{propext}, \texttt{Classical.choice},
\texttt{Quot.sound}, \texttt{finitePartition}).

\texttt{\textcolor{leancolor}{theorem} \textcolor{leanyellow}{changeSet\_constant} \{$\alpha$ : Type\} (d : TVal $\alpha$) \{t1 t2 : Time\}}

\texttt{~~(h : t1.val $\preceq$ t2.val)}

\texttt{~~(hno : $\lnot$ $\exists$ $\tau$ $\in$ changeSet d, $\tau$.val $\in$ Set.Ioc t1.val t2.val) :}

\texttt{~~interpValAt d t1 = interpValAt d t2 := by $\ldots$}
\label{changeSet_constant}
\end{restatable}

Therefore, if $\mathbf{I}\lsem d,\tau\rsem$ is defined for all $d\in \mathit{TVal}$ and all $\tau\in X_\tau$, by \ref{eq:df-nochange} it is defined for all $\tau\in\mathbf{T}$.

\comment{Do I really need to state this as a theorem and prove it?}

\subsection{Without features or time}

Considering the example above, choose \texttt{OurBikes} as $E$.  
Interpretation finds a ZFC-element of the ZFC-class for $E$.

\textbf{1.~~Create an instantiated tuple for $E$ at time $\tau=0$.}

Create a new triple with a unique identifier (UI):
\begin{gather}
\langle~\textbf{classifier},\texttt{"OurBikes\textbackslash\#a"}, \{x~\vert ~(~ x \in \texttt{OurBikes}~\wedge~\texttt{OurBikes}\in V_C~ \wedge \notag \\
 ~~~  \texttt{OurBikes}=\texttt{MyBike}\cup\texttt{YourBike}~\wedge~\#\texttt{OurBikes}=2~)~\}~\rangle
\end{gather}

\textbf{2.~~Recursively create instantiated triples for ZFC-classes used in the rule to define the class.}

Create instantiated triples for \texttt{MyBike} and (later) \texttt{YourBike}.
\begin{gather}
\langle~\textbf{classifier},\texttt{"MyBike\textbackslash\#b"},\{x~\vert ~(~ x \in \texttt{MyBike}~ \wedge \notag \\
 ~~~\texttt{MyBike}\in V_C~\wedge~ \texttt{MyBike}\subseteq\texttt{Bicycle}~\wedge \#\texttt{MyBike}=1~)~\}~\rangle
\end{gather}

\texttt{MyBike} uses \texttt{Bicycle} in its ZFC-class definition, so create an instantiated triple for \texttt{Bicycle}.
\begin{gather}
\langle~\textbf{classifier},\texttt{"Bicycle\textbackslash\#c"},\{x~\vert ~(~ x \in \texttt{Bicycle} ~\wedge~ \texttt{Bicycle}\in V_C~\wedge~ \texttt{Bicycle}\subseteq\texttt{Vehicle}~)~\}~\rangle
\end{gather}

\texttt{Bicycle} uses \texttt{Vehicle} in its ZFC-class definition, so create an instantiated triple for \texttt{Vehicle}.
\begin{gather}
\langle~\textbf{classifier},\texttt{"Vehicle\textbackslash\#d"},\{x~\vert ~(~ x \in \texttt{Vehicle} \wedge \texttt{Vehicle}\in V_C~)~\}~\rangle
\end{gather}

\texttt{Vehicle} uses no other ZFC-classes so is atomic.

Atomic triples are \bemph{complete}.  The UI of instantiated triples of atomic ZFC-classes represents a ZFC-element of that class.

\textbf{3.  Propagate the UI of a complete triple, back to the triple that referenced them, replacing the class name with the UI.}
\begin{gather}
\langle~\textbf{classifier},\texttt{"Bicycle\textbackslash\#c"},\{x~\vert ~(~ x \in \texttt{Bicycle} ~\wedge~ \notag \\
~~~ \texttt{Bicycle}\in V_C~\wedge~ \texttt{Bicycle}\subseteq\texttt{"Vehicle\textbackslash\#d"}~)~\}~\rangle
\end{gather}

\textbf{4.~~When all classes have been replaced by UIs, find an element of the class.}
\begin{gather}
\langle~\textbf{classifier},\texttt{"Bicycle\textbackslash\#c"},\texttt{"Vehicle\textbackslash\#d"}~\rangle
\end{gather}

Propagating \texttt{"Bicycle\textbackslash\#c"} to \texttt{"MyBike\textbackslash\#b"}:
\begin{gather}
\langle~\textbf{classifier},\texttt{"MyBike\textbackslash\#b"},\{x~\vert ~(~ x \in \texttt{MyBike}~ \wedge \notag \\
 ~~~\texttt{MyBike}\in V_C~\wedge~ \texttt{MyBike}\subseteq\texttt{"Bicycle\textbackslash\#c"}~\wedge \#\texttt{MyBike}=1~)~\}~\rangle
\end{gather}

and solving:
\begin{gather}
\langle~\textbf{classifier},\texttt{"MyBike\textbackslash\#b"},\texttt{"Bicycle\textbackslash\#c"}~\rangle
\end{gather}

Propagating \texttt{"MyBike\textbackslash\#b"} to \texttt{"OurBikes\textbackslash\#a"}:
\begin{gather}
\langle~\textbf{classifier},\texttt{"OurBikes\textbackslash\#a"}, \{x~\vert ~(~ x \in \texttt{OurBikes}~\wedge~\texttt{OurBikes}\in V_C~ \wedge \notag \\
 ~~~  \texttt{OurBikes}=\texttt{"MyBike\textbackslash\#b"}\cup\texttt{YourBike}~\wedge~\#\texttt{OurBikes}=2~)~\}~
\end{gather}

Then instantiate a triple for \texttt{YourBike}, which makes new triples for \texttt{Bicycle} and \texttt{Vehicle}, and solves them similarly to \texttt{MyBike}.

\begin{gather}
\langle~\textbf{classifier},\texttt{"YourBike\textbackslash\#e"},\texttt{"Bicycle\textbackslash\#f"}~\rangle \notag \\
\langle~\textbf{classifier},\texttt{"Bicycle\textbackslash\#f"},\texttt{"Vehicle\textbackslash\#g"}~\rangle \notag \\
\langle~\textbf{classifier},\texttt{"Vehicle\textbackslash\#g"},V_C~\rangle \notag \\
\end{gather}

Propagating \texttt{"YourBike\textbackslash\#e"} to \texttt{"OurBikes\textbackslash\#a"}:
\begin{gather}
\langle~\textbf{classifier},\texttt{"OurBikes\textbackslash\#a"}, \{x~\vert ~(~ x \in \texttt{OurBikes}~\wedge~\texttt{OurBikes}\in V_C~ \wedge \notag \\
 ~~~  \texttt{OurBikes}=\texttt{"MyBike\textbackslash\#b"}\cup\texttt{"YourBike\textbackslash\#e"}~\wedge~\#\texttt{OurBikes}=2~)~\}~\rangle
\end{gather}

Finally, \texttt{"OurBikes\textbackslash\#a"} has two ZFC-elements, but only one is needed.  (If this was a feature instead of a classifier, the result would be a pair (tuple) holding as many items as needed.)
\begin{gather}
\langle~\textbf{classifier},\texttt{"OurBikes\textbackslash\#a"}, \texttt{"YourBike\textbackslash\#e"}~\rangle
\end{gather}

So the tuples instated for $\mathbf{I}\lsem E,\tau\rsem$ are:
\begin{gather}
\langle~\textbf{classifier},\texttt{"OurBikes\textbackslash\#a"}, \texttt{"YourBike\textbackslash\#e"}~\rangle \notag \\
\langle~\textbf{classifier},\texttt{"MyBike\textbackslash\#b"},\texttt{"Bicycle\textbackslash\#c"}~\rangle \notag \\
\langle~\textbf{classifier},\texttt{"Bicycle\textbackslash\#c"},\texttt{"Vehicle\textbackslash\#d"}~\rangle \notag \\
\langle~\textbf{classifier},\texttt{"Vehicle\textbackslash\#d"},V_C~\rangle \notag \\
\langle~\textbf{classifier},\texttt{"YourBike\textbackslash\#e"},\texttt{"Bicycle\textbackslash\#f"}~\rangle \notag \\
\langle~\textbf{classifier},\texttt{"Bicycle\textbackslash\#f"},\texttt{"Vehicle\textbackslash\#g"}~\rangle \notag \\
\langle~\textbf{classifier},\texttt{"Vehicle\textbackslash\#g"},V_C~\rangle \notag \\
\end{gather}

The ZFC-element for $E$ in class \texttt{"OurBikes"} is \texttt{"OurBikes\textbackslash\#a"}, which by transitivity is also \texttt{"Bicycle\textbackslash\#f"} and \texttt{"Vehicle\textbackslash\#g"}.

Writing all this steps out is verbose, and won't be done again--just the important steps.
In practice it won't be done either, but this is how to find an element of the class represented by $d\in D$ to make $\mathbf{I}\lsem d,\tau\rsem$, to get 
$\mathbf{I}\lsem E,\tau\rsem$ for selected element $E\in D$.

\subsection{With features without time}

Take the example below to illustrate the procedure, a (non-association) classifier without connectors (features typed by associations).
\begin{lstlisting}
classifier Bicycle {
  feature rollsOn [2] : Wheel; 
  feature holdsWheel [*] : BikeFork;
}
classifier Wheel; 
classifier BikeFork;
\end{lstlisting}

The instantiation procedure starts with:
\textbf{1.~~Create a (design) tuple for each classifier or type}

\texttt{Wheel} and \texttt{BikeFork} have default membership in $V_C$:
\begin{gather}
\langle~\textbf{classifier},\texttt{"Wheel"},\{x~\vert ~(~ x \in \texttt{Wheel} \wedge \texttt{Wheel}\in V_C~)~\}~\rangle \notag \\
\langle~\textbf{classifier},\texttt{"BikeFork"},\{x~\vert ~(~ x \in \texttt{BikeFork} \wedge \texttt{BikeFork}\in V_C~)~\}~\rangle
\label{eq:wheel}
\end{gather}

\texttt{Bicycle} has two features, but the multiplicity of \texttt{holdsWheel} is undefined, \texttt{[*]}.  Presumably, some specialization of \texttt{Bicycle} would set the multiplicity of 
\texttt{holdsWheel}, or be determined by a consistency rule.  To make a model, when multiplicity is (still) a range, or undefined, the lowest allowed multiplicity will be arbitrarily selected.  Multiplicity of \texttt{[0]} is selected for \texttt{holdsWheel}, so is not instantiated as a feature of \texttt{Bicycle}.

By \ref{eq:df-typemultiplicity}, \texttt{classifier Bicycle \{ feature rollsOn [2] : Wheel; \}} 

\begin{gather}
\{~\langle \mathbf{classifier},w,C\rangle~\vert~w=W(\texttt{"Bicycle"})~\wedge~\{y_1,y_2\}\in\texttt{"Wheel"}~\wedge \notag \\
~~~C~=~\{\langle \texttt{"rollsOn"},R\rangle~\vert~R=\{~\langle w,y_1\rangle
\end{gather}

Then the Instantiation Procedure is very similar to that used for \texttt{OurBikes}.


